\documentclass{beamer}

\usepackage[T1]{fontenc}
\usepackage[utf8]{inputenc}
\usepackage{lmodern}
\usepackage{amsmath,amssymb,amsfonts}
\usepackage{bm}
\usepackage{physics}
\usepackage{mathtools}

\title{Transformers in Machine Learning}
\subtitle{Mathematical Structure and Connections to Physics and PDEs}
\author{}
\date{}

\begin{document}

%------------------------------------------------
\frame{\titlepage}

%------------------------------------------------
\begin{frame}{Outline}
\begin{enumerate}
\item Deep learning architectures in context
\item What is a transformer?
\item The mathematics of self-attention
\item Comparison with MLPs, CNNs, and RNNs
\item Multi-head attention and transformer blocks
\item Physics and statistical mechanics viewpoints
\item Transformers for PDEs and scientific machine learning
\item Summary
\end{enumerate}
\end{frame}

%------------------------------------------------
\begin{frame}{Deep Learning in Context}
Classical deep learning architectures include:
\begin{itemize}
\item multilayer perceptrons (MLPs),
\item convolutional neural networks (CNNs),
\item recurrent neural networks (RNNs).
\end{itemize}

Each architecture encodes a specific inductive bias:
\begin{itemize}
\item MLPs: dense global mixing,
\item CNNs: locality and translation structure,
\item RNNs: sequential recurrence and memory.
\end{itemize}

Transformers are also deep neural networks, but with a different structural principle:
\[
\textbf{adaptive interaction through attention.}
\]
\end{frame}
%------------------------------------------------
\begin{frame}{What Is a Transformer?}
  A transformer is a neural-network architecture built around the idea of
  \textbf{self-attention}.

  Core principle:
  \begin{itemize}
  \item each input element can interact directly with every other input element,
  \item the interaction strengths are learned and depend on the data itself.
  \end{itemize}

  Thus transformers replace fixed interaction patterns by
  \[
  \text{data-dependent couplings}.
  \]

  This makes them especially powerful for:
  \begin{itemize}
  \item language,
  \item long sequences,
  \item images,
  \item scientific fields and operator learning.
  \end{itemize}
\end{frame}

%------------------------------------------------
\begin{frame}{Input as a Sequence of Vectors}
  Suppose the input is a sequence
  \[
  x_1, x_2, \dots, x_n,
  \qquad x_i \in \mathbb{R}^d.
  \]

  Examples:
  \begin{itemize}
  \item words in a sentence,
  \item patches in an image,
  \item nodes in a discretized field,
  \item spatial or temporal samples of a PDE solution.
  \end{itemize}

  A transformer maps this sequence to another sequence of vectors:
  \[
  (x_1,\dots,x_n) \longmapsto (y_1,\dots,y_n).
  \]
\end{frame}

%------------------------------------------------
\begin{frame}{Self-Attention: The Basic Formula}
  For each position \(i\), the transformer forms
  \[
  y_i = \sum_{j=1}^n \alpha_{ij} v_j,
  \]
  where:
  \begin{itemize}
  \item \(v_j = W_V x_j\) are the \emph{values},
  \item \(\alpha_{ij}\) are attention weights.
  \end{itemize}

  The weights are computed from
  \[
  q_i = W_Q x_i, \qquad k_j = W_K x_j,
  \]
  called the \emph{query} and \emph{key} vectors.

  Then
  \[
  \alpha_{ij}
  =
  \frac{\exp\!\left(q_i \cdot k_j / \sqrt{d_k}\right)}
       {\sum_{\ell=1}^n \exp\!\left(q_i \cdot k_\ell / \sqrt{d_k}\right)}.
       \]

       This is the central mechanism of a transformer.
\end{frame}

%------------------------------------------------
\begin{frame}{Matrix Form of Attention}
  Collect the input vectors into a matrix
  \[
  X \in \mathbb{R}^{n \times d}.
  \]

  Define
  \[
  Q = X W_Q,\qquad K = X W_K,\qquad V = X W_V.
  \]

  Then the attention operation is
  \[
  \mathrm{Attention}(Q,K,V)
  =
  \mathrm{softmax}\!\left(\frac{QK^\top}{\sqrt{d_k}}\right)V.
  \]

  Here:
  \begin{itemize}
  \item \(QK^\top\) is an \(n\times n\) interaction matrix,
  \item the softmax normalizes each row,
  \item multiplication by \(V\) aggregates information.
  \end{itemize}
\end{frame}

%------------------------------------------------
\begin{frame}{Interpretation of the Attention Matrix}
  The matrix
  \[
  A = \mathrm{softmax}\!\left(\frac{QK^\top}{\sqrt{d_k}}\right)
  \]
  acts like a learned kernel or coupling matrix.

  Its entries satisfy
  \[
  A_{ij} \ge 0,
  \qquad
  \sum_j A_{ij}=1.
  \]

  Thus
  \[
  y_i = \sum_j A_{ij} v_j
  \]
  is a weighted average of the value vectors.

  Interpretation:
  \begin{itemize}
  \item query \(q_i\): ``what does position \(i\) look for?''
  \item key \(k_j\): ``what does position \(j\) offer?''
  \item value \(v_j\): ``what information does position \(j\) contribute?''
  \end{itemize}
\end{frame}

%------------------------------------------------
\begin{frame}{Comparison with a Standard MLP}
  A standard MLP layer has the form
  \[
  y = \sigma(Wx + b),
  \]
  with fixed weights \(W\).

  In contrast, attention uses
  \[
  y_i = \sum_j A_{ij}(X)\, v_j,
  \]
  where the effective coupling \(A_{ij}\) depends on the input \(X\).

  Thus:
  \[
  \text{MLP: fixed couplings}
  \qquad\text{vs}\qquad
  \text{Transformer: adaptive couplings}.
  \]

  This is one reason transformers are so expressive.
\end{frame}

%------------------------------------------------
\begin{frame}{Comparison with CNNs}
  In a convolutional neural network,
  \[
  y_i = \sum_{r \in \mathcal{N}(i)} w_r\, x_{i+r},
  \]
  where \(\mathcal{N}(i)\) is a small local neighborhood.

  Thus CNNs assume:
  \begin{itemize}
  \item locality,
  \item translation invariance,
  \item fixed kernels.
  \end{itemize}

  Attention instead uses
  \[
  y_i = \sum_j A_{ij}(X)\, x_j,
  \]
  which is
  \begin{itemize}
  \item global,
  \item adaptive,
  \item not restricted to local neighborhoods.
  \end{itemize}
\end{frame}

%------------------------------------------------
\begin{frame}{Comparison with RNNs}
  In an RNN,
  \[
  h_t = f(h_{t-1}, x_t),
  \]
  so information propagates sequentially.

  Advantages:
  \begin{itemize}
  \item natural for time series,
  \item explicit recurrence.
  \end{itemize}

  Limitations:
  \begin{itemize}
  \item long-range dependencies are hard,
  \item training can be unstable,
  \item computation is hard to parallelize.
  \end{itemize}

  Transformers avoid recurrence and instead connect all positions simultaneously through attention.
\end{frame}

%------------------------------------------------
\begin{frame}{Why the Factor \(1/\sqrt{d_k}\)?}
  The attention logits are
  \[
  q_i \cdot k_j.
  \]

  If \(q_i\) and \(k_j\) have \(d_k\) components with comparable variance, then typically
  \[
  q_i \cdot k_j \sim O(\sqrt{d_k})
  \]
  or \(O(d_k)\) depending on scaling assumptions.

  Without normalization, these logits can become large, causing the softmax to saturate.

  Therefore one rescales by
  \[
  \frac{1}{\sqrt{d_k}}
  \]
  to stabilize optimization and keep gradients in a useful regime.
\end{frame}

%------------------------------------------------
\begin{frame}{Multi-Head Attention}
  In practice, transformers use multiple attention heads.

  For head \(h\):
  \[
  Q^{(h)} = XW_Q^{(h)},\quad
  K^{(h)} = XW_K^{(h)},\quad
  V^{(h)} = XW_V^{(h)}.
  \]

  Each head computes
  \[
  \mathrm{head}_h
  =
  \mathrm{Attention}(Q^{(h)},K^{(h)},V^{(h)}).
  \]

  Then all heads are concatenated and linearly mixed:
  \[
  \mathrm{MultiHead}(X)
  =
  \mathrm{Concat}(\mathrm{head}_1,\dots,\mathrm{head}_H)W_O.
  \]

  Intuition:
  \begin{itemize}
  \item different heads learn different interaction patterns,
  \item one head may capture local structure, another long-range dependence.
  \end{itemize}
\end{frame}

%------------------------------------------------
\begin{frame}{The Transformer Block}
  A standard transformer block contains:

  \begin{enumerate}
  \item multi-head attention,
  \item residual connection,
  \item layer normalization,
  \item position-wise feedforward network,
  \item another residual connection and normalization.
  \end{enumerate}

  Schematically:
  \[
  X \mapsto X + \mathrm{MHA}(X)
  \]
  followed by
  \[
  X \mapsto X + \mathrm{MLP}(X).
  \]

  So a transformer is still a deep network built from familiar components:
  \[
  \text{attention} + \text{MLP} + \text{residual structure}.
  \]
\end{frame}

%------------------------------------------------
\begin{frame}{Positional Information}
  Attention itself is permutation-equivariant:
  \[
  (x_1,\dots,x_n)\mapsto (y_1,\dots,y_n)
  \]
  depends only on pairwise relations, not on absolute order.

  For sequence tasks we must add positional information.

  Common approaches:
  \begin{itemize}
  \item learned positional embeddings,
  \item sinusoidal encodings,
  \item relative position encodings.
  \end{itemize}

  For example, sinusoidal encodings use
  \[
  \mathrm{PE}(m,2r)=\sin\!\left(\frac{m}{10000^{2r/d}}\right),
  \qquad
  \mathrm{PE}(m,2r+1)=\cos\!\left(\frac{m}{10000^{2r/d}}\right).
  \]
\end{frame}

%------------------------------------------------
\begin{frame}{Transformers as Kernel Machines}
  Attention has the form
  \[
  y_i = \sum_j A_{ij}(X)\, v_j.
  \]

  This resembles an integral kernel operator:
  \[
  (\mathcal{K}f)(x)
  =
  \int K(x,x') f(x')\, dx'.
  \]

  In the discrete setting, attention acts like
  \[
  y_i = \sum_j K(x_i,x_j)\, v_j,
  \]
  but with a kernel \(K\) learned adaptively from the data.

  This viewpoint is useful in scientific machine learning and operator learning.
\end{frame}

%------------------------------------------------
\begin{frame}{Physics Viewpoint I: Adaptive Many-Body Couplings}
  A useful physics analogy is
  \[
  y_i = \sum_j J_{ij}(X)\, x_j,
  \]
  where
  \[
  J_{ij}(X) \equiv A_{ij}(X)
  \]
  is an input-dependent coupling.

  This resembles:
  \begin{itemize}
  \item mean-field interaction matrices,
  \item adaptive spin couplings,
  \item message-passing in many-body systems,
  \item self-consistent effective interactions.
  \end{itemize}

  So transformers can be interpreted as systems with
  \[
  \text{state-dependent interaction kernels}.
  \]
\end{frame}

%------------------------------------------------
\begin{frame}{Physics Viewpoint II: Statistical Mechanics Analogy}
  The softmax weights are
  \[
  A_{ij} =
  \frac{e^{s_{ij}}}{\sum_\ell e^{s_{i\ell}}},
  \qquad
  s_{ij} = \frac{q_i\cdot k_j}{\sqrt{d_k}}.
  \]

  This looks like a Gibbs or Boltzmann weight:
  \[
  p_j = \frac{e^{-\beta E_j}}{Z}.
  \]

  Indeed, if we identify
  \[
  E_{ij} = - s_{ij},
  \]
  then attention weights are Gibbs-like probabilities over interaction partners.

  This suggests a statistical-mechanics interpretation:
  \begin{itemize}
  \item scores \(s_{ij}\) define effective energies,
  \item softmax performs a local partition-function normalization.
  \end{itemize}
\end{frame}

%------------------------------------------------
\begin{frame}{Mean-Field Interpretation}
  Suppose each degree of freedom \(i\) updates by coupling to all others through effective coefficients \(A_{ij}\):
  \[
  y_i = \sum_j A_{ij} v_j.
  \]

  This resembles a mean-field update:
  \[
  m_i^{\text{new}} = F\!\left(\sum_j J_{ij} m_j\right),
  \]
  except that the couplings themselves depend on the current state.

  Thus transformers may be viewed as
  \[
  \text{nonlinear, adaptive mean-field models}.
  \]
\end{frame}

%------------------------------------------------
\begin{frame}{Attention and Graph Structure}
  If the input is viewed as a graph with nodes \(i\), then attention defines a complete graph with weighted edges
  \[
  i \longleftrightarrow j
  \]
  of strength \(A_{ij}\).

  This makes transformers closely related to:
  \begin{itemize}
  \item graph neural networks,
  \item message-passing networks,
  \item nonlocal interaction models.
  \end{itemize}

  Difference:
  \begin{itemize}
  \item graph neural networks often use a fixed graph,
  \item transformers learn the graph dynamically from the data.
  \end{itemize}
\end{frame}

%------------------------------------------------
\begin{frame}{Transformers and PDEs: Why They Matter}
  Many PDE problems involve long-range dependencies:
  \begin{itemize}
  \item elliptic equations,
  \item nonlocal operators,
  \item multiscale dynamics,
  \item global constraints.
  \end{itemize}

  CNNs are excellent for local structure, but transformers can naturally represent
  \[
  \text{nonlocal coupling across the whole domain}.
  \]

  This is one reason transformers have become important in scientific machine learning.
\end{frame}

%------------------------------------------------
\begin{frame}{Operator Learning Viewpoint}
  In PDEs one often wants to learn not just one solution, but an operator:
  \[
  \mathcal{G}: u_0(x) \mapsto u(x,t)
  \]
  or
  \[
  a(x) \mapsto u(x),
  \]
  where \(a(x)\) might be a coefficient field.

  Transformers can be used to approximate such mappings by treating the discretized field values as tokens and learning nonlocal interactions between them.

  This links transformers to:
  \begin{itemize}
  \item neural operators,
  \item Fourier neural operators,
  \item attention-based operator learning.
  \end{itemize}
\end{frame}

%------------------------------------------------
\begin{frame}{Example: Discretized PDE Fields as Tokens}
  Suppose a field \(u(x)\) is sampled at points
  \[
  x_1,\dots,x_n.
  \]

  Then each token may be
  \[
  x_i^{\text{token}} = \big(x_i, u(x_i), a(x_i), \dots \big).
  \]

  Attention then computes
  \[
  y_i = \sum_j A_{ij} v_j,
  \]
  which allows the representation at point \(x_i\) to depend on all other sampled points.

  This is particularly useful for:
  \begin{itemize}
  \item nonlocal PDEs,
  \item integral equations,
  \item strongly coupled multiscale systems.
  \end{itemize}
\end{frame}

%------------------------------------------------
\begin{frame}{Transformers vs Fourier Methods for PDEs}
  In scientific machine learning, Fourier-based methods and transformers both capture nonlocal structure.

  Fourier-based methods:
  \begin{itemize}
  \item use spectral decomposition,
  \item are natural for translation-invariant problems,
  \item often efficient on regular grids.
  \end{itemize}

  Transformers:
  \begin{itemize}
  \item are more flexible,
  \item handle irregular domains and adaptive relations,
  \item learn data-dependent kernels rather than fixed spectral ones.
  \end{itemize}
\end{frame}

%------------------------------------------------
\begin{frame}{Complexity of Self-Attention}
  A standard attention layer computes an \(n\times n\) matrix:
  \[
  QK^\top,
  \]
  so the cost scales as
  \[
  O(n^2 d).
  \]

  This is one of the main computational bottlenecks.

  By contrast:
  \begin{itemize}
  \item CNNs scale more locally,
  \item RNNs scale linearly in sequence length but sequentially in time.
  \end{itemize}

  This has motivated many sparse and efficient transformer variants.
\end{frame}

%------------------------------------------------
\begin{frame}{Why Transformers Became So Important}
  Transformers became dominant because they combine:
  \begin{itemize}
  \item global context,
  \item parallel computation,
  \item strong scaling with data and model size,
  \item flexible learned interaction structure.
  \end{itemize}

  They perform especially well when long-range correlations matter.

  In physics language:
  \[
  \text{they are powerful models for systems with rich nonlocal interactions.}
  \]
\end{frame}

%------------------------------------------------
\begin{frame}{Relation to Standard Deep Learning}
  Transformers are not separate from deep learning; they are a modern architecture within deep learning.

  They still use:
  \begin{itemize}
  \item linear layers,
  \item nonlinearities,
  \item residual connections,
  \item normalization,
  \item gradient-based training.
  \end{itemize}

  What is new is the central primitive:
  \[
  \boxed{\text{self-attention}}
  \]
  rather than convolution or recurrence.
\end{frame}

%------------------------------------------------
\begin{frame}{Strengths and Weaknesses}
  Strengths:
  \begin{itemize}
  \item captures long-range dependencies,
  \item highly parallelizable,
  \item adaptable to many data types.
  \end{itemize}

  Weaknesses:
  \begin{itemize}
  \item quadratic cost in sequence length,
  \item often data-hungry,
  \item can be less structured than domain-specific architectures.
  \end{itemize}

  In scientific applications, hybrids are often useful:
  \[
  \text{physics priors} + \text{attention mechanisms}.
  \]
\end{frame}

%------------------------------------------------
\begin{frame}{Transformers in Scientific Machine Learning}
  Active directions include:
  \begin{itemize}
  \item PDE surrogates,
  \item turbulence modeling,
  \item climate and weather forecasting,
  \item molecular and materials modeling,
  \item operator learning,
  \item quantum many-body representations.
  \end{itemize}

  The central attraction is that attention provides a flexible way to encode interactions across scales and distances.
\end{frame}

%------------------------------------------------
\begin{frame}{Summary}
  Transformers are deep neural networks built around self-attention.

  Mathematically:
  \[
  \mathrm{Attention}(Q,K,V)
  =
  \mathrm{softmax}\!\left(\frac{QK^\top}{\sqrt{d_k}}\right)V.
  \]

  Compared with standard methods:
  \begin{itemize}
  \item MLPs use fixed dense couplings,
  \item CNNs use local kernels,
  \item RNNs use sequential recurrence,
  \item transformers use adaptive global couplings.
  \end{itemize}

  This makes them especially powerful for nonlocal and multiscale problems.
\end{frame}

%------------------------------------------------
\begin{frame}{Physics Interpretation}
  A useful physics summary is:
  \begin{itemize}
  \item attention defines an input-dependent interaction matrix,
  \item softmax resembles Gibbs-type normalization,
  \item transformer layers resemble adaptive many-body updates,
  \item this connects naturally to statistical mechanics, graph models, and operator learning.
  \end{itemize}

  This is why transformers are becoming increasingly relevant in physics and PDE-based scientific machine learning.
\end{frame}

%------------------------------------------------
\begin{frame}{Possible Extensions}
  Natural follow-up topics include:
  \begin{itemize}
  \item transformers for quantum many-body states,
  \item graph transformers,
  \item neural operators and Fourier neural operators,
  \item attention in PDE solvers,
  \item statistical-mechanics interpretations of deep architectures.
  \end{itemize}
\end{frame}

\end{document}
